\section{Absolute differentials and the inseparable genus estimate}
\label{sec:differentials}

The two properties of the marked curve in \Cref{sec:twist} will now
give an exact lower bound for the arithmetic genus of its purely
inseparable covers. Throughout this section, a curve over a field $U$
means a regular integral projective scheme of dimension one over $U$.
It need not be geometrically regular or geometrically reduced. We set
\[
 \chi(C)=\dim_U H^0(C,\mathcal O_C)-\dim_U H^1(C,\mathcal O_C),
 \qquad \nu(C)=-2\chi(C).
\]
Divisor degrees and cohomology dimensions are always taken over the
stated ground field, and the degree of a vector bundle is the degree
of its determinant. Thus $\nu(C)$ retains any contribution from an
inseparable extension of the constant field.

Fix a finite group $J$ as in \Cref{thm:tuple-divisibility} and put
$m=|J|$. Only this order enters the following estimate, as a bound on
the cover degrees needed later.
Put $\kappa=\overline{\mathbb F}_p$. For each $s$, use the fields,
elliptic curve, and labels constructed in \Cref{prop:twist-index} for
that value of $s$. In particular,
$n=p^s$, $F'=\kappa(A_0^n)$, and $E_{F'}\simeq(A_0)_{F'}$.
Write $\mathcal B$ for the descended finite \'{e}tale divisor of labels;
over $F'$ it is the sum of the distinct rational points $t_1,\ldots,t_n$.
For $h\geq1$, $s\geq h$, and a finite extension $F/F_0$ with
$[F:F_0]_p\leq p^h$, we use two properties of the marked curve.
Every divisor degree on $E_F$ is divisible by $p^{s-h}$, by
\Cref{cor:index-base-change}. Over $S=FF'$, any set of the forms
$t_i^*\eta$ loses at most $h$ dimensions of independence, by
\Cref{lem:label-rank}, where $\eta$ is a nonzero invariant differential
on $A_0$.

\begin{theorem}[Exact purely inseparable genus estimate]
\label{thm:inseparable-genus}
Fix $h\ge1$ and take $s\ge h$ such that
\begin{equation}
 p^{s-h}>m\bigl(h+\log_p m\bigr).
 \label{eq:genus-rounding-threshold}
\end{equation}
Let $F/F_0$ be finite with $[F:F_0]_p\le p^h$, and let $Z$ be the
normalization of $E_F$ in a finite purely inseparable extension of its
function field of degree $d\le m$.  Choose a compositum $S=FF'$.
The curve $Z_S$ is integral and regular.  Let $r_i$ be the ramification
index of $Z_S\to E_S$ at the unique point over $t_i$.
With $\nu(Z)$ computed over $F$, one has
\begin{equation}
 \frac{\nu(Z)}d\ge\sum_{i=1}^n\left(1-\frac1{r_i}\right).
 \label{eq:inseparable-genus-bound}
\end{equation}
The bound is uniform in $F$, $Z$, and the tuple-class multiset used
later in the counting argument.
\end{theorem}

The proof has two parts. Absolute differentials give the same
inequality with an error of at most $h+\log_p d$. After multiplication
by $d$, the difference between the two sides is a difference of
divisor degrees over $F$, hence a multiple of $p^{s-h}$. The strict
threshold then forces this difference to be nonnegative. We develop
the differential identities first, together with the separable
ramification bound needed later in \Cref{sec:specialization}.

\subsection{Divisor degrees, duality, and separable ramification}

The required identities hold over more general fields than the
twisted-curve construction. Throughout this and the next subsection,
let $\kappa$ be any perfect field of characteristic $p$, let
$S/\kappa$ be finitely generated, and put
$\rho=\operatorname{trdeg}_\kappa S$. Curves and degrees in these two
subsections are over $S$.

The classical theory of inseparable genus change begins with Tate
\cite{Tate1952}; Schr\"oer \cite{Schroer2009} gives a modern account
for geometrically integral curves. Here absolute differentials retain
arithmetic genus without a geometric-reducedness hypothesis and also
measure ramification at the marked points.

\begin{lemma}[Degrees over the stated ground field]
\label{lem:curve-degrees}
For an invertible sheaf $\mathcal L$ on $C$,
\begin{equation}
 \chi(C,\mathcal L)=\deg_S\mathcal L+\chi(C).
 \label{eq:curve-euler-degree}
\end{equation}
A finite dominant morphism $f:C\to C'$ of such curves is finite flat.  If
its function-field degree is $\delta$, then
\[
 \deg_S f^*D=\delta\deg_S D
 \quad\text{and}\quad
 \deg_S f_*D'=\deg_S D'
\]
for divisors $D$ on $C'$ and $D'$ on $C$.  These assertions, and Euler
characteristics, add over disjoint unions of regular integral curves.
After a finite separable extension $S_1/S$, divisor and line-bundle degrees
and coherent Euler characteristics have the same numerical values over
$S_1$ as before extension over $S$.
\end{lemma}

\begin{proof}
For a closed point $x$ put $\deg_S(x)=[\kappa(x):S]$.  For any divisor $D$
the quotient in
\[
 0\longrightarrow\mathcal O_C(D)
 \longrightarrow\mathcal O_C(D+x)
 \longrightarrow\mathcal Q_x\longrightarrow0
\]
has length one over the DVR $\mathcal O_{C,x}$, so its dimension over $S$
is $[\kappa(x):S]$.  Finite-support sheaves have no higher cohomology.
Adding and subtracting closed points proves
$\chi(\mathcal O_C(D))=\chi(C)+\deg_S D$.  A rational trivialization of
$\mathcal L$ identifies it with $\mathcal O_C(D)$ for a divisor $D$, giving
\eqref{eq:curve-euler-degree}.  Applied to a principal divisor, the same
formula gives degree zero, so this degree depends only on $\mathcal L$.

Locally on $C'$, the finite module $f_*\mathcal O_C$ is torsion-free over
a DVR and therefore free, of generic rank $\delta$.
For a closed point $x\in C'$ and a uniformizer $t$ there, reduction of
this free module modulo $t$ gives
\begin{equation}
 \sum_{y\mapsto x}e_y[\kappa(y):\kappa(x)]=\delta,
 \label{eq:local-degree-formula}
\end{equation}
where $e_y$ is the valuation of $t$ in $\mathcal O_{C,y}$.
Indeed the summand at $y$ has a filtration of length $e_y$ with
successive quotients $\kappa(y)$.  Multiplication by
$[\kappa(x):S]$ proves the pullback degree formula.  The pushforward
formula follows from
$f_*y=[\kappa(y):\kappa(x)]x$ and the tower formula for residue degrees.

Finite separable scalar extension preserves regularity, and a regular
curve after such extension is a disjoint union of its integral
components.  Flat base change for coherent cohomology
in \Cref{foundation:cohomology} gives
\[
 H^i(C_{S_1},\mathcal G_{S_1})
   =H^i(C,\mathcal G)\otimes_S S_1.
\]
Thus the dimensions over the respective fields, and hence Euler
characteristics, agree.  Applying
\eqref{eq:curve-euler-degree} proves the assertion for line bundles and
their Cartier divisors.  In particular it also applies to a reduced
finite divisor: separable scalar extension preserves its reducedness.
\end{proof}

\begin{proposition}[The absolute determinant computes arithmetic genus]
\label{prop:absolute-genus}
The absolute cotangent sheaf $\Omega_{C/\kappa}$ is locally free of rank
$\rho+1$, and
\begin{equation}
 \deg_S\det\Omega_{C/\kappa}=\nu(C).
 \label{eq:absolute-determinant-genus}
\end{equation}
\end{proposition}

\begin{proof}
An affine coordinate ring of $C$ is a localization of a finitely
generated integral $\kappa$-algebra: choose a finitely generated domain
with fraction field $S$, clear the finitely many coefficients in a
presentation over $S$, and take the contraction of the defining prime
ideal before localizing.  Its local rings are consequently essentially
of finite type over $\kappa$.  The regularity criterion over a perfect
field in \Cref{foundation:regularity} says that these regular local rings
are local rings at smooth points of the corresponding
$\kappa$-variety.  Their absolute differentials are free of rank
\[
 \operatorname{trdeg}_\kappa\kappa(C)=\rho+1.
\]
The sheaf is coherent: the relative differentials over $S$ are coherent,
and $\Omega_{S/\kappa}$ is finite-dimensional, so the cotangent sequence
gives a finite presentation locally.  This proves the asserted local
freeness without making a smoothness assumption over $S$.

Choose a projective embedding $i:C\hookrightarrow\mathbb P^b_S$ and let
$\mathcal I$ be its ideal sheaf.  This is a regular immersion of
codimension $b-1$.  Here is a local verification.  Let $R$ be an ambient
regular local ring, $I$ the local ideal, and $R/I$ the regular local ring
of $C$.  Their dimensions differ by $b-1$, by the dimension formula at
the common residue field.  Since both are regular, the kernel of
\[
 \mathfrak m_R/\mathfrak m_R^2
   \longrightarrow
 \mathfrak m_{R/I}/\mathfrak m_{R/I}^2
\]
has dimension $b-1$.  Choose $b-1$ elements of $I$ lifting a basis of this
kernel.  They form part of a regular system of parameters of $R$ and
cut out a regular local quotient of dimension $\dim(R/I)$.
The remaining ideal defining $R/I$ in that quotient is prime and must
be zero: a nonzero prime in a finite-dimensional local domain strictly
lowers the quotient dimension.  Thus the chosen elements generate $I$.
It follows that $\mathcal I/\mathcal I^2$ is locally free of rank $b-1$.

The absolute conormal sequence is in fact short exact:
\begin{equation}
 0\longrightarrow\mathcal I/\mathcal I^2
 \longrightarrow i^*\Omega_{\mathbb P^b_S/\kappa}
 \longrightarrow\Omega_{C/\kappa}\longrightarrow0.
 \label{eq:absolute-conormal}
\end{equation}
To check the left end, the kernel of the last arrow is locally free,
since the last module is locally free.  Its rank is
$(b+\rho)-(1+\rho)=b-1$.  The conormal sequence surjects
$\mathcal I/\mathcal I^2$ onto this kernel.  A surjection between locally
free modules of the same rank is an isomorphism, proving
\eqref{eq:absolute-conormal}.

The projective space is obtained from $\mathbb P^b_\kappa$ by extending
constants.  Its absolute cotangents therefore split into relative
cotangents over $S$ and the constant bundle of $\Omega_{S/\kappa}$.
Taking determinants in \eqref{eq:absolute-conormal} gives
\[
 \det\Omega_{C/\kappa}
  \simeq\mathcal L_c\otimes_S\det\Omega_{S/\kappa},
 \qquad
 \mathcal L_c=
 \mathcal O_C(-b-1)\otimes
       \det(\mathcal I/\mathcal I^2)^\vee.
\]
The last tensor factor is a constant line and has degree zero.

We compute the degree of $\mathcal L_c$ using the
projective-space duality in \Cref{foundation:duality}.
A Koszul resolution for the local regular sequence defining $C$ shows
that
\[
 \mathcal E xt^v_{\mathbb P^b_S}
       (i_*\mathcal O_C,\mathcal O(-b-1))
 =\begin{cases}
   i_*\mathcal L_c,&v=b-1,\\
   0,&v\ne b-1.
  \end{cases}
\]
After dualizing the Koszul resolution, its top cohomology
is the dual of the top exterior power of the free module of equations,
restricted to $C$.  A change of those equations acts on this exterior
power by the determinant of their change on $\mathcal I/\mathcal I^2$.
Thus the local top cohomology modules glue to
$\det(\mathcal I/\mathcal I^2)^\vee\otimes\mathcal O_C(-b-1)$.

The local-to-global Ext spectral sequence has only this one nonzero
row.  Projective-space duality consequently gives, for $r=0,1$,
\[
 H^r(C,\mathcal O_C)^\vee
  \simeq
 \operatorname{Ext}^{b-r}_{\mathbb P^b_S}
       (i_*\mathcal O_C,\mathcal O(-b-1))
  \simeq H^{1-r}(C,\mathcal L_c).
\]
Hence $\chi(C,\mathcal L_c)=-\chi(C)$.  Applying
\eqref{eq:curve-euler-degree} yields
$\deg_S\mathcal L_c=-2\chi(C)$, which proves
\eqref{eq:absolute-determinant-genus}.
\end{proof}

\begin{lemma}[A separable determinant bound, including the wild term]
\label{lem:separable-different}
Let $f:C\to C'$ be a finite dominant, generically separable morphism of
regular integral projective curves over $S$, of degree $\delta$.
There is an effective divisor $R_f$ on $C$ such that
\[
 \nu(C)=\delta\nu(C')+\deg_S R_f.
\]
If $y\in C$ has ramification index $a_y$ over $f(y)$, then
\begin{equation}
 \operatorname{ord}_y R_f\ge a_y-1+\mathbf 1_{p\mid a_y}.
 \label{eq:separable-different-bound}
\end{equation}
No separability of $\kappa(y)/\kappa(f(y))$ is required.
The formula and the inequalities can be added componentwise after a
finite separable extension of $S$.
\end{lemma}

\begin{proof}
The morphism
$f^*\Omega_{C'/\kappa}\to\Omega_{C/\kappa}$ is generically an isomorphism:
at the function fields a finite separable extension has no relative
differentials, and both absolute differential spaces have dimension
$\rho+1$.  Its determinant is therefore a nonzero regular morphism of
line bundles.  The divisor of this morphism is effective; call it $R_f$.
Its degree is the difference of the two determinant degrees, which is
$\nu(C)-\delta\nu(C')$ by
\Cref{lem:curve-degrees,prop:absolute-genus}.

Write $x=f(y)$ and choose a uniformizer $t$ of the DVR
$B=\mathcal O_{C',x}$.  Its absolute differential is primitive in
$\Omega_{B/\kappa}$, meaning that it is part of a $B$-basis.  Indeed the
conormal sequence for the residue field is
\[
 (t)/(t^2)\longrightarrow
 \Omega_{B/\kappa}\otimes_B\kappa(x)
 \longrightarrow\Omega_{\kappa(x)/\kappa}\longrightarrow0.
\]
The middle term has dimension $\rho+1$ and the right term dimension
$\rho$, since $\kappa(x)$ is finite over $S$ and
\Cref{foundation:fields} computes absolute differential dimension over
the perfect field $\kappa$.  Thus the first map has nonzero
one-dimensional image, generated by $dt$.  Consequently $dt$ is
nonzero modulo the maximal ideal, and it extends to a basis of the
free module $\Omega_{B/\kappa}$.

In $A=\mathcal O_{C,y}$ write $t=u\tau^{a_y}$, where $u$ is a unit and
$\tau$ a uniformizer.  Absolute differentiation gives
\[
 dt=\tau^{a_y}\,du+
       a_yu\tau^{a_y-1}\,d\tau.
\]
All coordinates of this column of the differential matrix are
divisible by $\tau^{a_y-1}$; when $p\mid a_y$, the second term is zero,
so they are divisible by $\tau^{a_y}$.  Expanding the determinant in
this column proves \eqref{eq:separable-different-bound}.  Every step
uses absolute differentials over $\kappa$, so an inseparable extension
of residue fields causes no change to the argument.
\end{proof}

\subsection{The differential calculation for a height-one map}

We have established the determinant formula for arithmetic genus and
the separable ramification bound. The remaining ingredient is an
exact genus identity for a purely inseparable map of exponent one;
successive such maps will handle arbitrary purely inseparable covers.

Write $T_C=(\Omega_{C/\kappa})^\vee$ for the absolute tangent bundle.
It has rank $\rho+1$, even when the relative tangent sheaf over $S$ has
different behavior.

The calculation below is a curve-level form of the canonical-bundle
formula for a $p$-closed foliation; compare
\cite[Section~I.1, Proposition~1.1 and Equations~(1.2)--(1.3)]{Ekedahl1988}.
We prove the version needed here directly from local presentations,
using absolute differentials over the perfect constants. In this way
the kernel, the Frobenius-power cokernel, and all inseparable residue
extensions remain explicit.

\begin{lemma}[Height-one kernels and cokernels]
\label{lem:height-one}
Let $\pi:C\to C'$ be a finite dominant morphism of regular integral
projective curves over $S$ such that
$\kappa(C)^p\subseteq\kappa(C')$.  Put
$q=[\kappa(C):\kappa(C')]=p^a$ and
\[
 \mathcal F=\ker(T_C\longrightarrow\pi^*T_{C'}),
 \qquad
 \mathcal Q=\operatorname{coker}(T_C\longrightarrow\pi^*T_{C'}).
\]
Both $\mathcal F$ and $\mathcal Q$ are locally free of rank $a$, and
\begin{equation}
 \det\mathcal Q\simeq(\det\mathcal F)^{\otimes p}.
 \label{eq:height-one-cokernel}
\end{equation}
Consequently
\begin{equation}
 \nu(C)=q\nu(C')+(p-1)\deg_S\mathcal F.
 \label{eq:height-one-genus}
\end{equation}
\end{lemma}

\begin{proof}
\textit{The integral presentation.}
On an affine open of $C'$, every element of the finite upstairs
coordinate ring has its $p$th power in the downstairs ring: the power
lies in the downstairs fraction field and is integral, so normality
puts it in that ring.  There is therefore a unique prime upstairs over
each downstairs prime, since membership of an element is determined
by membership of its $p$th power.  At a closed point we may consequently
write $B\subset A$ for the downstairs and upstairs DVRs, with $A$ still
finite over $B$.  It is finite free of rank $q$ and satisfies
$A^p\subset B$.

Let $\tau$ be an upstairs uniformizer, let $\epsilon$ be the
ramification index, and put $f=[\kappa(A):\kappa(B)]$.
The element $\tau^p\in B$ has upstairs valuation $p$; all nonzero
elements of $B$ have upstairs valuation in $\epsilon\mathbb Z$.
Thus $\epsilon\mid p$, so $\epsilon$ is $1$ or $p$.
Reduction of the rank-$q$ free module modulo a downstairs uniformizer
gives $q=\epsilon f$ by \eqref{eq:local-degree-formula}.
The residue extension has exponent at most one.  Successively
adjoining a generator outside the preceding residue field gives a
degree-$p$ extension each time.  Choose residue generators
$\bar b_1,\ldots,\bar b_r$ in this manner, so their monomials with
exponents in $\{0,\ldots,p-1\}$ form a $\kappa(B)$-basis of $\kappa(A)$
and $f=p^r$.  Lift them to $b_1,\ldots,b_r\in A$.  If $\epsilon=p$,
append $b_{r+1}=\tau$; if $\epsilon=1$, append nothing.  In both cases
there are exactly $a$ generators, and $c_\ell=b_\ell^p$ belongs to $B$.

For a downstairs uniformizer $t$, the ideal $tA$ is
$\tau^\epsilon A$.  Filter $A/tA$ by the powers of $\tau$ from $0$ to
$\epsilon$.  Each successive quotient is $\kappa(A)$, and the lifted
residue monomials give its basis after multiplication by the
corresponding power of $\tau$.  It follows that the $q$ monomials
\[
 b_1^{i_1}\cdots b_a^{i_a},\qquad 0\le i_\ell<p,
\]
are a $\kappa(B)$-basis of $A/tA$.
By Nakayama they generate $A$ over $B$; since $A$ is free of rank $q$,
they are a $B$-basis.  The homomorphism from the monic-relation algebra
therefore gives the actual presentation
\begin{equation}
 A\simeq
 B[X_1,\ldots,X_a]/(X_1^p-c_1,\ldots,X_a^p-c_a),
 \qquad X_\ell\longmapsto b_\ell.
 \label{eq:height-one-dvr-presentation}
\end{equation}
Both sides have the same free monomial basis over $B$, which verifies
that there are no further relations.  This includes the cases of a
nontrivial inseparable residue extension and simultaneous residue
extension and ramification.

These presentations are available on open neighborhoods.  Shrink an
open downstairs around the chosen point so that all the $b_\ell$ are
sections of $\pi_*\mathcal O_C$, all the $c_\ell$ are regular downstairs,
and the determinant of the displayed monomial basis is invertible.
Then \eqref{eq:height-one-dvr-presentation} holds there.  Such
neighborhoods give local presentations throughout the curves.

\textit{Free kernel and free cokernel.}
Set $N=A\otimes_B\Omega_{B/\kappa}$, and let $W\subset N$ be the span
of $dc_1,\ldots,dc_a$.  Differentiating the displayed presentation,
using characteristic $p$, gives
\begin{equation}
 \Omega_{A/\kappa}\simeq
       (N/W)\oplus\bigoplus_{\ell=1}^a A\,db_\ell.
 \label{eq:height-one-differential-presentation}
\end{equation}
Both $N$ and $\Omega_{A/\kappa}$ are free of rank $\rho+1$ by
\Cref{prop:absolute-genus}.  The summand $N/W$ in
\eqref{eq:height-one-differential-presentation} is therefore free,
of rank $\rho+1-a$.  The sequence
$0\to W\to N\to N/W\to0$ splits, so $W$ is free of rank $a$.
Its $a$ generators $dc_\ell$ are consequently a basis.

Dualizing now shows that the tangent kernel consists of the
$\kappa$-derivations $A\to A$ killing $B$, and has free coordinates
\[
 D\longmapsto(D(b_1),\ldots,D(b_a)).
\]
The tangent cokernel is $W^\vee$, with coordinates
\[
 [\delta]\longmapsto
       (\delta(c_1),\ldots,\delta(c_a)),
 \qquad \delta\in\operatorname{Der}_\kappa(B,A).
\]
Indeed a derivation $B\to A$ extends to $A$ exactly when it kills
the $c_\ell$; when it does, the values assigned to the $b_\ell$ are
arbitrary.  Thus both kernel and cokernel are locally free of rank $a$.

\textit{Gluing the determinant lines.}
On an overlap, write another system of generators as
$b'_k=P_k(b_1,\ldots,b_a)$ with coefficients in $B$, and put
\[
 \mathsf J_{k\ell}=
 \frac{\partial P_k}{\partial X_\ell}(b_1,\ldots,b_a).
\]
A derivation in the kernel kills the coefficients, so its new
coordinate vector is $\mathsf J$ times its old coordinate vector.
The matrix $\mathsf J$ is invertible since both coordinate systems are bases.
Taking $p$th powers of the generator identities gives downstairs
\[
 c'_k=(b'_k)^p=P_k^{[p]}(c_1,\ldots,c_a),
\]
where $P_k^{[p]}$ means that the coefficients of $P_k$ are raised to
the $p$th power.  Their absolute differentials vanish.  Differentiating
this identity by a representative $\delta:B\to A$ of a cokernel class
therefore gives
\[
 \delta(c'_k)=\sum_\ell \mathsf J_{k\ell}^{p}\,\delta(c_\ell).
\]
Thus the cokernel coordinate transition is the entrywise Frobenius
power $\mathsf J^{[p]}$.  Coordinate transitions for determinant lines are
$\det\mathsf J$ and $\det(\mathsf J^{[p]})=(\det\mathsf J)^p$; if one uses local frames
instead, both transitions are inverted.  In either convention they
give exactly \eqref{eq:height-one-cokernel}, with the exponent $+p$.

Finally the exact sequence
\[
 0\longrightarrow\mathcal F\longrightarrow T_C
 \longrightarrow\pi^*T_{C'}\longrightarrow\mathcal Q
 \longrightarrow0
\]
and \eqref{eq:height-one-cokernel} give
\[
 \deg T_C=q\deg T_{C'}-(p-1)\deg\mathcal F.
\]
Since $\deg T_C=-\nu(C)$, this is \eqref{eq:height-one-genus}.
\end{proof}

\subsection{Proof of the exact estimate}

Return to the fields and marked curve in
\Cref{thm:inseparable-genus}, with $\kappa=\overline{\mathbb F}_p$
and $S=FF'$. The height-one identity will be applied along a
Frobenius tower from $Z_S$ to $E_S$. At each stage the label
covectors bound the degree of the tangent-map image; the resulting
weighted sum gives the bounded-error estimate before the final
index argument removes the error.

\begin{proof}[Proof of \Cref{thm:inseparable-genus}]
\textit{Normalization and the separable splitting field.}
Normalizations are finite by \Cref{foundation:normalization}, and normal
curves over fields are regular.  We first record why the finite purely
inseparable maps in this argument are radicial, that is, universally
injective.  If $A$ is the integral closure of a normal affine domain
$B$ in a finite purely inseparable extension, there is an exponent $v$
such that $a^{p^v}\in B$ for every $a\in A$.  The power belongs to
$\operatorname{Frac}(B)$ and is integral over $B$, so it belongs to $B$
by normality.  After tensoring with any $B$-algebra, the $p^v$th power
of every element still lies in the image of that algebra.  Consequently
over a prime of the base there is at most one prime upstairs: membership
of an element in a prime is determined by membership of its $p^v$th
power downstairs.  The residue extension is purely inseparable for
the same reason.  This proves the assertion after every base change.

The extension $S/F$ is finite separable, since $F'/F_0$ is finite
Galois.  Thus $Z_S$ is regular and reduced by
\Cref{foundation:regularity}.  Its map to the integral curve $E_S$ is
still finite, surjective, and radicial.  This map is a homeomorphism
on underlying spaces, so $Z_S$ is irreducible and hence integral.
It is the normalization of $E_S$ in its function field: it is finite
and normal with that field, and an element integral over the base is
also integral over its finite normal coordinate ring.  Its degree
remains $d$ by finite flatness.  By \Cref{lem:curve-degrees},
\[
 \nu(Z_S)=\nu(Z),
\]
where the two Euler characteristics are computed over their stated
fields.  We may therefore do the differential calculation over $S$.
If $d=1$, this curve is $E_S$, every $r_i=1$, and the theorem holds.
Assume henceforth that $d>1$.

\textit{The Frobenius tower and its recurrence.}
Let $M_e=S(E)$ and $M_z=S(Z)$, and set
\[
 M_j=M_eM_z^{p^j}\quad(0\le j\le b),\qquad
 M_0=M_z,\qquad M_b=M_e,
\]
where $b$ is the first index for which the last equality holds.
Let $Y_j$ be the normalization of $E_S$ in $M_j$, and write
\[
 \gamma_j:Y_j\longrightarrow E_S,\qquad
 d_j=[M_j:M_e],\qquad
 \pi_j:Y_j\longrightarrow Y_{j+1},\qquad
 q_j=d_j/d_{j+1}=p^{m_j}.
\]
In particular $Y_0=Z_S$, $d_0=d$, $Y_b=E_S$, and $d_b=1$.
The defining fields satisfy
\begin{equation}
 M_{j+1}=M_eM_j^p.
 \label{eq:frobenius-tower-field}
\end{equation}
Thus each $\pi_j$ has exponent at most one and
\Cref{lem:height-one} applies.

Put $T_j=(\Omega_{Y_j/\kappa})^\vee$ and
\[
 \mathcal F_j=\ker(T_j\longrightarrow\pi_j^*T_{j+1}).
\]
This is also the kernel of the composite tangent map
$T_j\to\gamma_j^*T_{E_S}$.  To see this, at the function field a
$\kappa$-derivation kills $M_e$ if and only if it kills
$M_eM_j^p=M_{j+1}$, since every derivation kills $p$th powers.
The two kernels are therefore equal generically.  They are equal as
sheaves: a section whose image vanishes generically has zero image
everywhere in either of the locally free, hence torsion-free, targets.

The genus identity involves $\deg_S\mathcal F_j$. We will estimate
this degree through the image of the tangent map to $E_S$.
Choose a nonzero invariant differential $\eta$ on $A_0$.
The constant-curve identification of $E_S$ gives the splitting
\[
 \Omega_{E_S/\kappa}
   \simeq\mathcal O_{E_S}\eta
      \oplus(\mathcal O_{E_S}\otimes_S\Omega_{S/\kappa}).
\]
Thus its dual is a trivial bundle of rank $n+1$. In this trivial
target we can use the marked points to impose linear vanishing
conditions on the tangent-map image.

Let $\mathcal H_j$ be the image in $\gamma_j^*T_{E_S}$, and put
$b_j^\circ=-\deg_S\mathcal H_j$.
It is a torsion-free coherent sheaf on a regular curve, hence locally
free.  Its rank is $n+1-m_j$, since
$\operatorname{trdeg}_\kappa S=n$.
The exact sequence with kernel $\mathcal F_j$ and image $\mathcal H_j$,
together with \eqref{eq:height-one-genus}, gives
\[
 \nu(Y_j)=q_j\nu(Y_{j+1})+(p-1)\deg_S\mathcal F_j,
 \qquad
 b_j^\circ=\nu(Y_j)+\deg_S\mathcal F_j.
\]
Eliminating $\deg_S\mathcal F_j$ and using $d_j=q_jd_{j+1}$ yields
\begin{equation}
 \frac{\nu(Y_j)}{d_j}
  =\frac1p\frac{\nu(Y_{j+1})}{d_{j+1}}
     +\left(1-\frac1p\right)\frac{b_j^\circ}{d_j}.
 \label{eq:genus-recurrence}
\end{equation}
The terminal term $\nu(Y_b)$ is zero because $E_S$ has genus one.
It remains to give a lower bound for
$b_j^\circ/d_j=-\deg_S\mathcal H_j/d_j$.

\textit{Vanishing at a ramified label.}
For each $j$ let $S_j$ be the set of labels at which $\gamma_j$ has
ramification index greater than one.  For $i\in S_j$ write
$D_{j,i}=\gamma_j^*(t_i)$ and
\[
 v_i=(1,-t_i^*\eta)\in S\oplus\Omega_{S/\kappa}.
\]
Regard $v_i$ as a constant linear functional on
$\gamma_j^*T_{E_S}$. We claim that its restriction to $\mathcal H_j$
vanishes along the entire divisor $D_{j,i}$; equivalently, its values
lie in $\mathcal O_{Y_j}(-D_{j,i})$.

Let $a$ be a local parameter defining the rational point $t_i$ on
$E_S$.  Restriction of absolute differentials along $t_i$ sends
$(\lambda,\beta)$ in the displayed splitting to
$\lambda t_i^*\eta+\beta$.  Its kernel is the line spanned by $v_i$.
Since $a$ restricts to zero, $da$ belongs to this kernel modulo $(a)$.
Moreover its relative component is nonzero: $t_i$ is an $S$-rational
point on a smooth curve over $S$, and $a$ generates its relative
cotangent space.  Hence
\[
 da\equiv\lambda_i v_i\pmod{a\Omega_{E_S/\kappa}}
 \quad\text{for some }\lambda_i\in S^\times.
\]
At the unique point upstairs write $a=u\tau^r$, with $u$ a unit.
The nontrivial purely inseparable index $r$ is divisible by $p$;
therefore
\[
 d(\gamma_j^*a)=\tau^r\,du=(\gamma_j^*a)\,u^{-1}du.
\]
It lies in the ideal of the \emph{full} divisor $D_{j,i}$ times
$\Omega_{Y_j/\kappa}$.  Pulling back the preceding congruence and
dividing by the unit $\lambda_i$ shows that the pullback of $v_i$ lies
in that same ideal times $\Omega_{Y_j/\kappa}$.
To relate this differential to the claimed functional, denote the
tangent map by $d\gamma_j:T_j\to\gamma_j^*T_{E_S}$ and write
$\gamma_j^*v_i$ for the image of the pulled-back covector in
$\Omega_{Y_j/\kappa}$. For a local section $\theta$ of $T_j$, duality
gives
\[
 v_i(d\gamma_j(\theta))=(\gamma_j^*v_i)(\theta).
\]
The divisibility just proved makes this value a section of
$\mathcal O_{Y_j}(-D_{j,i})$. Since $\mathcal H_j$ is the image of
$d\gamma_j$, the claim follows. The vanishing is along
$\gamma_j^*(t_i)$ with its full multiplicity, not just its reduced
support.

\textit{From independent vanishing conditions to a degree bound.}
By \Cref{lem:label-rank}, among the $v_i$ for $i\in S_j$ there are
$c\ge |S_j|-h$ independent covectors over $S$.

Complete the $c$ chosen covectors to a constant basis of the dual of
the target trivial tangent bundle.  The full-divisor constraint gives
an injection
\begin{equation}
 \mathcal H_j\hookrightarrow
 \mathcal O_{Y_j}^{n+1-c}\oplus
       \bigoplus_{i\ \mathrm{chosen}}\mathcal O_{Y_j}(-D_{j,i}).
 \label{eq:label-image-injection}
\end{equation}
Let $R=n+1-m_j$ be the rank of $\mathcal H_j$.
The $R$th exterior power of the bundle on the right is a direct sum
of line bundles, each the tensor product of a choice of $R$ summands.
The induced
map from $\det\mathcal H_j$ is nonzero, so it has a nonzero component
in at least one of these lines. Only $n+1-c$ of the available
summands are trivial. The chosen line must therefore contain at least
$R-(n+1-c)=c-m_j$ negative divisor summands when this number is
positive, and at least zero otherwise.
Each $D_{j,i}$ has degree $d_j$ by finite flatness and rationality of
$t_i$. If the nonzero component uses $k$ negative summands, its target
line has degree $-kd_j$.  A nonzero regular map from one line bundle
to another has an effective zero divisor, and consequently
\[
 \deg_S\mathcal H_j\le-kd_j,
 \qquad
 \frac{b_j^\circ}{d_j}\ge c-m_j\ge |S_j|-h-m_j.
\]
The argument also covers rank zero, with $\det(0)=\mathcal O$.

\textit{The weighted estimate.}
Write $r_i=p^{u_i}$.  Every height-one step has local index at most
$p$, as proved in \Cref{lem:height-one}. The composite
$Y_0\to Y_j$ consequently has ramification index at most $p^j$ at
this label. In particular, $u_i\leq b$. Since ramification indices
multiply, the remaining map
$Y_j\to E_S$ has index at least $p^{u_i-j}$ whenever $j<u_i$.
Hence $i\in S_j$ for every $0\le j<u_i$. Iterating
\eqref{eq:genus-recurrence} with terminal value zero gives
\[
 \frac{\nu(Z)}d
  =\sum_{j=0}^{b-1}
      p^{-j}\left(1-\frac1p\right)\frac{b_j^\circ}{d_j}.
\]
Write $w_j=p^{-j}(1-1/p)$ for the weight of stage $j$. For each
label, the weights from those first $u_i$ stages sum to
\[
 \sum_{j=0}^{u_i-1}p^{-j}\left(1-\frac1p\right)
      =1-p^{-u_i}=1-\frac1{r_i}.
\]
Summing these contributions by label gives
\[
 \sum_{j=0}^{b-1}w_j|S_j|
  =\sum_{i=1}^n\sum_{\substack{0\leq j<b\\i\in S_j}}w_j
  \geq\sum_{i=1}^n\left(1-\frac1{r_i}\right).
\]
The error terms in the degree bound contribute separately. Since
$m_j\le\log_p d$ and the sum of all the weights is $1-p^{-b}<1$,
\[
 \sum_{j=0}^{b-1}w_j(h+m_j)\leq h+\log_p d.
\]
Substitute $b_j^\circ/d_j\geq |S_j|-h-m_j$ in the recurrence sum.
Subtracting the error bound from the label contribution yields
\begin{equation}
 \frac{\nu(Z)}d
   \ge\sum_{i=1}^n\left(1-\frac1{r_i}\right)
          -\bigl(h+\log_p d\bigr).
 \label{eq:genus-approximate}
\end{equation}

\textit{Divisors over $F$ and exact rounding.}
Multiplying \eqref{eq:genus-approximate} by $d$ bounds the possible
deficit by $d(h+\log_p d)$. Set $N=p^{s-h}$ and consider the difference
\[
 A=\nu(Z)-d\sum_i\left(1-\frac1{r_i}\right).
\]
We will realize both terms as degrees of divisors on $Z$, each
retaining its degree under pushforward to $E_F$. The index property
\Cref{cor:index-base-change} then makes both degrees divisible by
$N$. These divisors must be defined over $F$, before the separable
extension to $S$.
By \Cref{prop:absolute-genus}, the integer $\nu(Z)$ is the degree over
$F$ of $\det\Omega_{Z/\kappa}$.  A rational trivialization represents
this line bundle by a Cartier divisor on the regular integral curve
$Z$.  Its pushforward to $E_F$ preserves degree over $F$, by
\Cref{lem:curve-degrees}.  Thus
\[
 \nu(Z)\in N\mathbb Z.
\]
This uses the absolute determinant on $Z/F$ and its Euler
characteristic over $F$, including any inseparable constant field.

Let $T$ be the reduced inverse image in $Z$ of $\mathcal B_F$.
It is an effective divisor on the regular curve $Z$.
Finite separable scalar extension preserves reducedness, so $T_S$
is the reduced inverse image of the split labels.  If $z_i$ is the
unique point above the rational point $t_i$, the local degree formula
is
\[
 r_i[\kappa(z_i):S]=d.
\]
It follows that
\[
 \deg_F T=\deg_S T_S=\sum_i\frac d{r_i}.
\]
In particular this sum counts the full residue degrees of the reduced
points; those residue fields may be inseparable over $S$.
The divisor $\gamma^*\mathcal B_F-T$ is defined on $Z$ over $F$, where
$\gamma:Z\to E_F$, and its degree is
\[
 dn-\sum_i\frac d{r_i}
    =d\sum_i\left(1-\frac1{r_i}\right).
\]
Push it to $E_F$ to conclude that this integer also belongs to
$N\mathbb Z$.

Consequently $A$ is a multiple of $N$. By
\eqref{eq:genus-approximate} and $d\le m$,
\[
 A\ge-d(h+\log_p d)
      \ge-m(h+\log_p m)>-N.
\]
A multiple of $N$ strictly greater than $-N$ is nonnegative.  This
proves \eqref{eq:inseparable-genus-bound}.
The sufficient threshold \eqref{eq:genus-rounding-threshold} contains
no datum from $F$, $Z$, or a tuple-class multiset.  If a threshold
independent of the prime is desired, the stronger condition
\[
 2^{s-h}>m(h+\log_2m)
\]
suffices for every $p\ge2$.  This proves the stated uniformity.
\end{proof}
